\documentclass[../DoAn.tex]{subfiles}
\begin{document}

Chương~\ref{chapter:Introduction} đã giới thiệu bài toán. Chương này đối chiếu tài liệu quy trình nghiệp vụ do nhóm xây dựng với kịch bản sử dụng và các nhóm công cụ quản lý lịch phổ biến, sau đó xác định tác nhân, use case và quy tắc cần triển khai. Các nhận định về nhu cầu được dùng làm giả thiết thiết kế, chưa phải kết quả khảo sát định lượng tại một đơn vị cụ thể.

\section{Khảo sát hiện trạng}
\label{section:2.1}

\subsection{Khảo sát người dùng và quy trình hiện hành}
Kịch bản sử dụng của nhóm đặt ra yêu cầu người mượn cần tra cứu đồng thời phòng trống theo khung giờ, sức chứa, trang thiết bị và vị trí trên tầng. Mô hình tám tầng và sơ đồ chữ U là dữ liệu minh họa cho ứng dụng. Các yêu cầu này được kiểm tra qua luồng nghiệp vụ và giao diện, chưa được trình bày như kết quả phỏng vấn người dùng thực tế.

Nguồn nghiệp vụ chính là tài liệu \emph{Quy trình nghiệp vụ RoomHub} do nhóm xây dựng \cite{quytrinh2026}. Tài liệu mô tả các bước từ tra cứu, gửi yêu cầu, phê duyệt, tiếp cận phòng tới hoàn tất. Hai phương thức tiếp cận được xét là nhận khóa cơ hoặc thẻ từ và tạo mã tạm thời cho khóa số quản lý trực tuyến. Bảng~\ref{tab:loaiphong} tóm tắt các trường hợp minh họa trong quy trình.

\begin{table}[H]
\centering\small
\renewcommand{\arraystretch}{1.2}
\caption{Các cấu hình phòng trong kịch bản nghiệp vụ}
\label{tab:loaiphong}
\begin{tabularx}{\textwidth}{|p{1.6cm}|p{4.2cm}|X|}
\hline
\textbf{Loại} & \textbf{Phương thức} & \textbf{Quy trình tiếp cận} \\ \hline
Loại A & Chỉ dùng mã số & Nhân viên cơ sở vật chất đặt mã tại khóa trước giờ sử dụng và thay đổi sau khi kết thúc. \\ \hline
Loại B & Chỉ dùng thẻ từ & Người dùng nhận thẻ tại điểm bàn giao, sử dụng và trả thẻ sau khi kết thúc. \\ \hline
Loại C & Có cả mã số và thẻ; phiên sử dụng chọn mã & Áp dụng quy trình cấp mã, không mặc nhiên bàn giao thẻ. \\ \hline
Loại D & Có cả mã số và thẻ; phiên sử dụng chọn thẻ & Áp dụng quy trình nhận và trả thẻ, không mặc nhiên cấp mã. \\ \hline
\end{tabularx}
\end{table}

Trong phần mô phỏng khóa số, nhóm dùng cấu hình SmartLock kết nối nền tảng OneIoT. Ứng dụng khái quát phương thức tiếp cận thành khóa cơ/thẻ từ và khóa mã số trực tuyến. Mỗi yêu cầu dùng một phương thức theo loại khóa của phòng; quyền tạo mã được gắn với đúng phòng sau khi yêu cầu được duyệt.

\subsection{Khảo sát các giải pháp tương tự}
Để định vị giải pháp, nhóm đối chiếu RoomHub với các nhóm công cụ thường dùng để ghi lịch và đặt chỗ. Bảng~\ref{tab:sosanh} là so sánh định tính theo khả năng thiết kế, không phải phép đo hay đánh giá toàn bộ sản phẩm trên thị trường.

\begin{table}[H]
\centering\small
\renewcommand{\arraystretch}{1.2}
\caption{So sánh các nhóm giải pháp quản lý đặt phòng}
\label{tab:sosanh}
\begin{tabularx}{\textwidth}{|p{3.6cm}|X|X|X|X|}
\hline
\textbf{Tiêu chí} & \textbf{Sổ ghi / bảng tính} & \textbf{Lịch tài nguyên trực tuyến} & \textbf{Phần mềm đặt chỗ thương mại} & \textbf{RoomHub} \\ \hline
Trạng thái phòng theo khung giờ & Thủ công & Có & Có & Có \\ \hline
Sơ đồ tầng, vị trí phòng & Không & Không & Thường có & Có (chữ U, 8 tầng) \\ \hline
Lọc sức chứa, thiết bị & Không & Hạn chế & Có & Có \\ \hline
Quy tắc riêng (hạn đặt trước, giới hạn yêu cầu, bảo trì) & Không & Hạn chế & Có, cấu hình theo gói & Có, quản trị viên chỉnh trong ứng dụng \\ \hline
Quy trình phê duyệt & Thủ công & Tự động chấp nhận hoặc từ chối & Có & Có, kèm đổi phòng \\ \hline
Cấp quyền tiếp cận theo loại khóa & Không & Không & Qua tích hợp riêng & Hẹn nhận khóa / mã PIN tạm thời \\ \hline
Chi phí, khả năng tùy biến & Công sức quản lý thủ công & Phụ thuộc gói dịch vụ & Phụ thuộc nhà cung cấp & Tự phát triển \\ \hline
\end{tabularx}
\end{table}

Từ phân tích trên, nhóm ưu tiên sơ đồ tầng có trạng thái theo thời gian, thao tác đặt phòng ngay sau khi chọn phòng, kiểm tra quy tắc trước khi lưu, quy trình duyệt và đổi phòng có điều kiện, quyền tiếp cận theo loại khóa, cùng khả năng chỉnh tham số nghiệp vụ trong ứng dụng.

\section{Tổng quan chức năng}
\label{section:2.2}
Mục này tóm tắt chức năng của hệ thống ở mức tổng quan; đặc tả chi tiết được trình bày trong mục~\ref{section:2.3}.

\subsection{Biểu đồ use case tổng quát}
\label{subsection:2.2.1}
Hệ thống có hai vai trò người dùng chính: giảng viên gửi và theo dõi yêu cầu, cán bộ quản lý duyệt và cấu hình. OneIoT và khóa thông minh là hệ thống ngoài trong nhánh cấp mã. Hình~\ref{fig:uc_tongquat} thể hiện use case ở mức khái niệm; phần truyền bản tin hiện được ứng dụng Android thực hiện trực tiếp. Vai trò được xác định từ tài khoản đăng nhập.

\begin{figure}[H]
    \centering
    \includegraphics[width=0.95\textwidth]{Hinhve/uc_tongquat.pdf}
    \caption{Biểu đồ use case tổng quát của hệ thống RoomHub}
    \label{fig:uc_tongquat}
\end{figure}

Các use case dùng chung gồm đăng nhập, đăng ký, khôi phục mật khẩu, xem thông báo và cập nhật hồ sơ. Người dùng có các chức năng tìm và đặt phòng, xem lịch theo trạng thái, quản lý yêu cầu của mình, ủy quyền người nhận khóa hộ và tạo mã PIN tạm thời khi được cấp quyền. Quản trị viên quản lý tài khoản, phòng, yêu cầu, lịch, cấu hình và khóa; thao tác duyệt, từ chối hoặc đổi phòng là phần trung tâm của quy trình.

\subsection{Biểu đồ use case phân rã ``Tìm và đặt phòng''}
\label{subsection:2.2.2}
Use case ``Tìm và đặt phòng'' được phân rã như Hình~\ref{fig:uc_datphong}. Người dùng bắt buộc chọn khung thời gian, chọn tầng, xem trạng thái phòng trên sơ đồ, chọn phòng và gửi yêu cầu; việc lọc theo sức chứa, thiết bị và loại khóa là tùy chọn. Khi gửi, hệ thống luôn thực hiện use case ``Kiểm tra quy tắc nghiệp vụ'' với các quy tắc trong Bảng~\ref{tab:br}.

\begin{figure}[H]
    \centering
    \includegraphics[width=0.9\textwidth]{Hinhve/uc_datphong.pdf}
    \caption{Biểu đồ use case phân rã ``Tìm và đặt phòng''}
    \label{fig:uc_datphong}
\end{figure}

\subsection{Biểu đồ use case phân rã ``Duyệt / từ chối / đổi phòng''}
\label{subsection:2.2.3}
Hình~\ref{fig:uc_quanly} phân rã use case quản lý yêu cầu đặt phòng của quản trị viên. Quản trị viên xem các yêu cầu chờ, sau đó có thể duyệt, từ chối, đổi sang phòng phù hợp, hẹn lịch nhận khóa cơ hoặc thẻ, tạo mã PIN hoặc thu hồi quyền tự tạo mã. Khi duyệt yêu cầu tại phòng khóa số, hệ thống cấp quyền tự tạo mã cho đúng người đặt và phòng đó; quyền này không thay đổi vai trò tài khoản.

\begin{figure}[H]
    \centering
    \includegraphics[width=0.95\textwidth]{Hinhve/uc_quanly.pdf}
    \caption{Biểu đồ use case phân rã ``Duyệt / từ chối / đổi phòng''}
    \label{fig:uc_quanly}
\end{figure}

\subsection{Quy trình nghiệp vụ}
\label{subsection:2.2.4}
Quy trình đặt phòng kết hợp thao tác của hai vai trò như Hình~\ref{fig:hoatdong}. Người dùng chọn thời gian và phòng, nhập mục đích rồi gửi yêu cầu. Sau khi duyệt, phòng khóa số có thể tạo mã tạm thời; ứng dụng lưu mã trước và gửi bản tin qua OneIoT khi có kết nối. Với khóa cơ hoặc thẻ từ, hai bên thống nhất thời điểm nhận khóa và có thể ủy quyền người nhận hộ. Biểu đồ phản ánh luồng nghiệp vụ tổng quát; trạng thái phát bản tin không đồng nghĩa khóa đã nhận mã.

\begin{figure}[H]
    \centering
    \includegraphics[width=\textwidth,height=0.84\textheight,keepaspectratio]{Hinhve/hoatdong.pdf}
    \caption{Biểu đồ hoạt động của quy trình đặt phòng và cấp quyền tiếp cận}
    \label{fig:hoatdong}
\end{figure}

Vòng đời yêu cầu được mô tả ở Hình~\ref{fig:trangthai}. Bản triển khai lưu các trạng thái chờ duyệt, đã duyệt, từ chối, đã hủy và vắng mặt; thông tin check-in/check-out cùng phương thức tiếp cận được quản lý bằng các trường hoặc đối tượng liên quan. Trạng thái mã tạm thời hiển thị theo thời gian và dấu thu hồi. Hình là sơ đồ khái niệm, nên không đồng nhất mọi bước sử dụng phòng với một trạng thái lưu trữ riêng.

\begin{figure}[H]
    \centering
    \includegraphics[width=0.95\textwidth]{Hinhve/trangthai.pdf}
\caption{Biểu đồ trạng thái yêu cầu đặt phòng và mã PIN tạm thời}
    \label{fig:trangthai}
\end{figure}

\section{Đặc tả chức năng}
\label{section:2.3}
Mục này đặc tả năm use case quan trọng nhất. Các use case còn lại gồm hủy đặt phòng, hẹn lịch nhận khóa và ủy quyền, quản lý lịch bảo trì, cấu hình nghiệp vụ và quản lý SmartLock được đặc tả tại Phụ lục~\ref{appendix:usecase}. Trước hết, Bảng~\ref{tab:br} liệt kê các quy tắc nghiệp vụ mà hệ thống đã hiện thực; các quy tắc được đánh mã để tham chiếu trong đặc tả và kiểm thử.

\begin{table}[H]
\centering\small
\renewcommand{\arraystretch}{1.15}
\caption{Các quy tắc nghiệp vụ đã hiện thực trong RoomHub}
\label{tab:br}
\begin{tabularx}{\textwidth}{|p{1.3cm}|X|p{3.4cm}|}
\hline
\textbf{Mã} & \textbf{Quy tắc} & \textbf{Tham số mặc định} \\ \hline
BR-01 & Phòng phải tồn tại và ở trạng thái có thể sử dụng. & -- \\ \hline
BR-02 & Ngày và giờ phải hợp lệ; giờ bắt đầu sớm hơn giờ kết thúc. & -- \\ \hline
BR-03 & Chỉ được đặt trong khoảng ngày cấu hình; có thể cho phép cùng ngày nếu giờ kết thúc chưa qua. & Mặc định 1 đến 3 ngày \\ \hline
BR-04 & Không được đặt phòng trùng với lịch bảo trì chưa bị hủy của phòng đó. & -- \\ \hline
BR-05 & Không trùng với yêu cầu đang chờ hoặc đã duyệt khác của cùng phòng; một người không có hai yêu cầu chồng thời gian. & -- \\ \hline
BR-06 & Tổng số yêu cầu đang chờ hoặc sắp sử dụng của một người không vượt giới hạn; mục đích sử dụng không được để trống. & 2 yêu cầu \\ \hline
BR-07 & Chỉ duyệt yêu cầu chưa hết hiệu lực; kiểm tra lại bảo trì và trùng lịch. & -- \\ \hline
BR-08 & Người đặt chỉ hủy yêu cầu đã duyệt trước giờ bắt đầu một khoảng tối thiểu. & 30 phút \\ \hline
BR-09 & Quản trị viên chỉ đổi phòng trước giờ bắt đầu một khoảng tối thiểu; phòng thay thế phải cùng loại khóa, đủ sức chứa và thiết bị, không bảo trì và không trùng lịch. & 30 phút \\ \hline
BR-10 & Mã PIN chỉ tạo cho yêu cầu đã duyệt tại phòng khóa số đang gắn khóa thông minh; hiệu lực có khoảng đệm quanh giờ sử dụng và bị thu hồi khi hủy hoặc đổi phòng. & 10 phút \\ \hline
BR-11 & Người đặt chỉ tự tạo mã khi có quyền còn hiệu lực cho đúng phòng; quản trị viên có thể thu hồi quyền theo phòng. & -- \\ \hline
BR-12 & Ủy quyền nhận hộ và hẹn nhận khóa chỉ áp dụng cho yêu cầu đã duyệt tại phòng dùng khóa cơ hoặc thẻ; lịch hẹn phải trước giờ sử dụng. & -- \\ \hline
\end{tabularx}
\end{table}

\subsection{Đặc tả use case ``Đăng nhập''}
\begin{dacta}{Đặc tả use case UC01 -- Đăng nhập}{tab:uc01}
Mã use case & UC01 \\ \hline
Tác nhân & Người dùng, quản trị viên \\ \hline
Mô tả & Người dùng xác thực bằng tên đăng nhập và mật khẩu; hệ thống xác định vai trò và mở giao diện tương ứng. \\ \hline
Tiền điều kiện & Tài khoản đã tồn tại và đang ở trạng thái hoạt động. \\ \hline
Luồng chính &
\begin{buoc}
\item Người dùng mở ứng dụng, nhập tài khoản và mật khẩu, bấm ``Đăng nhập''.
\item Hệ thống chuẩn hóa tên đăng nhập (bỏ khoảng trắng, chuyển chữ thường).
\item Hệ thống xác thực tài khoản đang hoạt động; ở chế độ máy chủ, dịch vụ Java tạo phiên đăng nhập.
\item Hệ thống mở ``Chế độ quản trị viên'' hoặc ``Chế độ người dùng'' theo vai trò tài khoản.
\end{buoc} \\ \hline
Luồng phát sinh & 3a. Sai thông tin hoặc tài khoản bị khóa: hiển thị lỗi và giữ nguyên màn hình. 3b. Người dùng chọn ``Đăng ký'' hoặc ``Quên mật khẩu?'' để tạo tài khoản mới hoặc đặt lại mật khẩu bằng mã khôi phục. \\ \hline
Hậu điều kiện & Phiên đăng nhập được thiết lập; mọi yêu cầu tiếp theo tới máy chủ mang token Bearer. \\
\end{dacta}

\subsection{Đặc tả use case ``Tìm và đặt phòng''}
\begin{dacta}{Đặc tả use case UC02 -- Tìm và đặt phòng}{tab:uc02}
Mã use case & UC02 \\ \hline
Tác nhân & Người dùng \\ \hline
Tiền điều kiện & Người dùng đã đăng nhập và tài khoản còn hoạt động. \\ \hline
Luồng chính &
\begin{buoc}
\item Người dùng mở ``Tìm và đặt phòng''; hệ thống sinh danh sách ngày được phép theo cấu hình (mặc định từ ngày mai đến ba ngày sau).
\item Người dùng chọn ngày, giờ bắt đầu, giờ kết thúc và tầng từ 1 đến 8.
\item Người dùng có thể lọc theo sức chứa tối thiểu, thiết bị cần có và loại khóa.
\item Hệ thống hiển thị bảy phòng của tầng trên sơ đồ chữ U, mỗi ô ghi sức chứa, thiết bị và tô màu trống, đã đặt, bảo trì hoặc không khớp lọc; đồng thời báo số phòng phù hợp đang trống.
\item Người dùng chọn một phòng trống (hoặc nhập tên phòng để tìm), nhập mục đích và gửi yêu cầu đặt phòng.
\item Hệ thống kiểm tra BR-01 đến BR-06, lưu yêu cầu chờ duyệt, thông báo cho quản trị viên và làm mới sơ đồ.
\end{buoc} \\ \hline
Luồng phát sinh & 5a. Phòng được chọn không ở trạng thái trống: yêu cầu chọn phòng khác. 6a. Vi phạm quy tắc: hiển thị lý do và không lưu yêu cầu. \\ \hline
Hậu điều kiện & Yêu cầu chờ duyệt được lưu; phòng hiển thị là đã có yêu cầu trong khung giờ đó; quản trị viên nhận thông báo. \\
\end{dacta}

\subsection{Đặc tả use case ``Duyệt yêu cầu đặt phòng''}
\begin{dacta}{Đặc tả use case UC03 -- Duyệt yêu cầu đặt phòng}{tab:uc03}
Mã use case & UC03 \\ \hline
Tác nhân & Quản trị viên \\ \hline
Tiền điều kiện & Quản trị viên đã đăng nhập; có yêu cầu đang chờ duyệt. \\ \hline
Luồng chính &
\begin{buoc}
\item Quản trị viên mở danh sách yêu cầu đặt phòng hoặc thông báo về yêu cầu mới.
\item Hệ thống hiển thị thông tin người đặt, phòng, ngày giờ, mục đích.
\item Quản trị viên chọn duyệt.
\item Hệ thống khóa bản ghi yêu cầu, kiểm tra BR-07: yêu cầu còn chờ, chưa qua giờ kết thúc, không trùng bảo trì hay yêu cầu đã duyệt khác.
\item Hệ thống chuyển sang trạng thái đã duyệt, ghi người và thời điểm xử lý.
\item Nếu phòng dùng khóa mã số, hệ thống cấp quyền tự tạo mã cho người đặt tại phòng đó.
\item Hệ thống gửi thông báo ``Yêu cầu đã được duyệt'' cho người đặt.
\end{buoc} \\ \hline
Luồng phát sinh & 3a. Quản trị viên chọn từ chối: hệ thống lưu kết quả và thông báo cho người đặt. 4a. Vi phạm BR-07: hiển thị lý do, trạng thái không đổi. \\ \hline
Hậu điều kiện & Yêu cầu được duyệt hoặc bị từ chối; người đặt nhận thông báo. Với phòng khóa số, quyền tự tạo mã theo phòng được cấp khi duyệt. \\
\end{dacta}

\subsection{Đặc tả use case ``Đổi phòng thay thế''}
\begin{dacta}{Đặc tả use case UC04 -- Đổi phòng thay thế}{tab:uc04}
Mã use case & UC04 \\ \hline
Tác nhân & Quản trị viên \\ \hline
Tiền điều kiện & Yêu cầu đã được duyệt và còn ít nhất 30 phút (theo cấu hình) trước giờ bắt đầu. \\ \hline
Luồng chính &
\begin{buoc}
\item Quản trị viên chọn yêu cầu cần đổi, ví dụ khi phòng cũ phải phục vụ sự kiện ưu tiên.
\item Hệ thống chỉ đề xuất phòng khác đang khả dụng, cùng loại khóa, có sức chứa và thiết bị không kém phòng cũ.
\item Quản trị viên chọn phòng thay thế và nhập lý do.
\item Hệ thống kiểm tra BR-09 cùng lịch bảo trì và các yêu cầu chờ/đã duyệt của phòng mới.
\item Hệ thống cập nhật phòng, ghi lịch sử đổi phòng, thu hồi mã PIN cũ và lịch hẹn nhận khóa cũ; nếu phòng mới dùng khóa số thì cấp quyền tạo mã cho phòng mới.
\item Hệ thống thông báo cho người đặt kèm lý do đổi phòng.
\end{buoc} \\ \hline
Luồng phát sinh & 4a. Phòng thay thế khác loại khóa: báo ``Phòng thay thế phải có cùng loại khóa với phòng hiện tại''. 4b. Quá hạn đổi phòng hoặc thiếu lý do: báo lỗi tương ứng. \\ \hline
Hậu điều kiện & Yêu cầu được chuyển sang phòng mới và ghi lịch sử đổi phòng; quyền tiếp cận phòng cũ bị vô hiệu. \\
\end{dacta}

\subsection{Đặc tả use case ``Tạo mã PIN tạm thời''}
\begin{dacta}{Đặc tả use case UC05 -- Tạo mã PIN tạm thời}{tab:uc05}
Mã use case & UC05 \\ \hline
Tác nhân & Người đặt hoặc quản trị viên; nền tảng OneIoT \\ \hline
Tiền điều kiện & Yêu cầu đã duyệt tại phòng khóa số đang gắn khóa thông minh; người đặt có quyền tự tạo mã theo phòng. Có thể tạo mã trong ứng dụng khi chưa kết nối OneIoT. \\ \hline
Luồng chính &
\begin{buoc}
\item Người đặt mở danh sách yêu cầu của mình, chọn yêu cầu và tạo mã.
\item Hệ thống kiểm tra BR-10, BR-11 và bảo đảm yêu cầu chưa có mã đang hiệu lực.
\item Hệ thống sinh mã bảy chữ số, tính thời điểm hiệu lực từ giờ bắt đầu trừ 10 phút đến giờ kết thúc cộng 10 phút theo cấu hình mặc định.
\item Hệ thống cấp vị trí mật khẩu của khóa và lưu mã cùng khoảng hiệu lực vào yêu cầu.
\item Nếu OneIoT đang kết nối, ứng dụng mã hóa và phát bản tin tạo mã qua MQTT; nếu chưa kết nối, mã chờ gửi trong ứng dụng.
\item Hệ thống hiển thị mã và thông báo rõ trạng thái đã gửi hoặc chờ kết nối.
\end{buoc} \\ \hline
Luồng phát sinh & 2a. Quyền đã bị thu hồi: báo ``Quyền tự tạo mã tại phòng này đã bị thu hồi hoặc chưa được cấp''. 4a. Phòng chưa gắn khóa: báo ``Phòng chưa được gắn với SmartLock đang quản lý''. 5a. Chưa kết nối hoặc gửi lỗi: mã vẫn được lưu để thử gửi khi kết nối lại; mã hết hiệu lực hoặc đã thu hồi sẽ không gửi. \\ \hline
Hậu điều kiện & Ứng dụng lưu mã và trạng thái gửi; việc khóa thật áp dụng mã chỉ được xác nhận bằng phản hồi hoặc kiểm thử thiết bị. \\
\end{dacta}

\section{Yêu cầu phi chức năng}
\label{section:2.4}
Về \textbf{tính đúng đắn và độ tin cậy}, hệ thống kiểm tra xung đột theo khoảng thời gian nửa mở, lịch bảo trì và quyền của người gọi. Ở chế độ máy chủ, thao tác ghi được kiểm tra trong giao dịch có khóa; tính đúng đắn khi nhiều thiết bị truy cập đồng thời cần được tiếp tục kiểm thử với máy chủ triển khai thực tế.

Về \textbf{bảo mật}, các API cần phiên đăng nhập hợp lệ theo từng loại thao tác. Token OneIoT được nhập cho phiên kết nối và không nằm trong kho mã nguồn. Ứng dụng mã hóa mật khẩu trước khi phát bản tin MQTT; phần lưu trữ mã tại ứng dụng và máy chủ cần được gia cố trước triển khai thực tế.

Về \textbf{tính khả dụng}, giao diện tiếng Việt và điều khiển chọn ngày giờ phù hợp thao tác cảm ứng trên điện thoại Android. Bản Release phải khởi động độc lập với Metro; chế độ dùng máy chủ còn phụ thuộc kết nối tới dịch vụ đang chạy.

Về \textbf{tính dễ bảo trì và mở rộng}, mã nguồn được chia theo module nghiệp vụ, các tham số chính sách được lưu dưới dạng cấu hình thay vì hằng số trong giao diện, và việc chuyển giữa lưu cục bộ với đồng bộ máy chủ chỉ phụ thuộc một cờ cấu hình.

Về \textbf{yêu cầu kỹ thuật của bản triển khai}, máy chủ dùng Java 17, Spring Boot, JPA/Hibernate và H2 trong môi trường thử nghiệm; ứng dụng dùng React Native, TypeScript, Async Storage và module Android viết bằng Kotlin để giao tiếp OneIoT. Đây là lựa chọn triển khai của nhóm, không phải tập công nghệ được suy diễn là bắt buộc của mọi học phần.

Chương này đã xác định vai trò, các use case, quy tắc nghiệp vụ và yêu cầu phi chức năng. Các yêu cầu về giao dịch, phân quyền, dữ liệu cục bộ và bản tin khóa là cơ sở cho lựa chọn công nghệ ở Chương~\ref{chapter:Methodology}.

\end{document}
